\documentclass[12pt]{article}

\usepackage{amsmath,amssymb,hyperref}
\topmargin -.5cm
\textheight 21cm
\oddsidemargin -.125cm
\textwidth 17cm

\newcommand{\D}{\Delta}
\newcommand{\cC}{\mathcal{C}}
\newcommand{\cO}{\mathcal{O}}

\begin{document}
\sloppy

\begin{center}
{\Large\bf Axion Frequency from the \texorpdfstring{$SO(2,4)$}{SO(2,4)} Casimir}
\end{center}

\section{Purpose}

This note sketches how the axion frequency should be obtained if one imitates the logic of
\texttt{2504.12290}.  The point is to separate three distinct ingredients:
\begin{itemize}
\item the local linearized equation in the projective $X^a$ patch,
\item the global covariance condition under $SO(2,4)$,
\item the extra physical condition that the fluctuation be the normalizable scalar primary.
\end{itemize}

The main claim is simple:
\begin{equation}
\text{Casimir} + \text{primary condition} + \text{normalizability}
\quad \Longrightarrow \quad
\omega=4
\end{equation}
for the massless RR axion zero mode on $AdS_5\times S^5$.

However, this note is about the \emph{strict} large-$R$ plus \texttt{2504.12290} logic.  In that
framework one should not solve the exact radial equation on the full spacetime and then read off
$\omega$.  Instead one should use:
\begin{itemize}
\item the local large-$R$ linearized equation in one projective patch,
\item the global covariance conditions supplied by the extended-BV/isometry formalism,
\item and the identification of the resulting global representation.
\end{itemize}
So the real question is not ``can the local radial equation be solved exactly?'' but rather
``how much of $\omega$ is fixed by local large-$R$ dynamics plus global representation theory?''

\section{What the Local Patch Does and Does Not Do}

For the axion zero mode on the sphere, the projective-coordinate ansatz is
\begin{equation}
\phi(X)=e^{-i\omega\tau(X)}f(s),
\qquad
s=\frac{r^2}{R^2},
\qquad
r^2=\sum_{m=1}^4 (X^m)^2.
\end{equation}
The exact scalar equation gives the radial ODE
\begin{equation}
4s(1+s)f''(s)+4(2+3s)f'(s)+\frac{\omega^2}{1+s}f(s)=0.
\label{radialode}
\end{equation}
Expanding around the center,
\begin{equation}
f(s)=\sum_{n=0}^{\infty} a_n s^n,
\end{equation}
one gets a regular local solution for every chosen $\omega$.  So the inner weak-curvature expansion does
\emph{not} quantize $\omega$.

This is exactly the same structural issue emphasized in \texttt{2504.12290}: one perturbative patch knows
only about a neighborhood of the chosen origin, not about the full target space.

\section{The \texttt{2504.12290} Logic}

The point of \texttt{2504.12290} is that the fluctuation problem should be formulated globally, not only in
one local patch.  Schematically, a fluctuation string field $\Phi$ is required to satisfy:
\begin{equation}
\text{on-shell equation}
\qquad\text{and}\qquad
\text{covariance under the global isometries.}
\end{equation}
For the compact example of that paper, the fluctuations are organized into definite isometry
representations, and applying the covariance condition twice identifies the target-space Laplacian with the
quadratic Casimir of the isometry group.

The AdS$_5$ axion analogue should have exactly the same structure.  The relevant global symmetry is now
$SO(2,4)$ rather than a compact isometry group.  Since the axion zero mode on $S^5$ is a scalar both on
AdS and on the sphere, no extra spin or sphere-Casimir terms are present.

So the fluctuation problem should reduce schematically to
\begin{equation}
\Box_{AdS}\phi=0,
\qquad
\frac12 M_{AB}M^{AB}\phi=\cC_2\,\phi,
\label{schematic}
\end{equation}
with the geometric identification
\begin{equation}
\frac12 M_{AB}M^{AB}\ \leftrightarrow\ \Box_{AdS}
\end{equation}
on the scalar fluctuation.

\section{The Scalar \texorpdfstring{$SO(2,4)$}{SO(2,4)} Casimir}

For a scalar primary of dimension $\D$, the relevant lowest-weight representation is
\begin{equation}
D(\D;0,0).
\end{equation}
Its quadratic Casimir is
\begin{equation}
\cC_2\big(D(\D;0,0)\big)=\D(\D-4).
\label{casimir}
\end{equation}

For the massless axion zero mode, the five-dimensional mass is zero, so the AdS wave equation gives
\begin{equation}
\Box_{AdS}\phi=0.
\end{equation}
Using \eqref{schematic}, this means
\begin{equation}
\cC_2=0.
\end{equation}
Therefore \eqref{casimir} implies
\begin{equation}
\D(\D-4)=0.
\label{dbranches}
\end{equation}
So there are two algebraic possibilities:
\begin{equation}
\D=0,\qquad \D=4.
\end{equation}

\section{What the Strict Large-\texorpdfstring{$R$}{R} Argument Actually Gives}

In the strict large-$R$ framework, the local projective-patch analysis gives only the inner recursion
problem.  One writes
\begin{equation}
\phi(X)=e^{-i\omega\tau(X)}f(s),
\qquad
s=\frac{r^2}{R^2},
\end{equation}
and the linearized EOM imply the exact radial ODE \eqref{radialode}.  Expanding for small $s$,
\begin{equation}
f(s)=a_0+a_1 s+a_2 s^2+\cdots,
\end{equation}
determines all higher coefficients in terms of $a_0$ and $\omega$.  Thus the local large-$R$ analysis
produces a one-parameter family of regular local solutions.  It does \emph{not} quantize $\omega$.

What the \texttt{2504.12290} logic adds is not an exact radial solve, but a global representation
condition.  Once the fluctuation is promoted from a local patch solution to a globally covariant object,
one may ask which $SO(2,4)$ representation it belongs to.  For the axion, the dynamical input from the
linearized EOM says that the relevant fluctuation sector is a massless scalar.  Therefore the global
representation must be one of the scalar lowest-weight modules with
\begin{equation}
\cC_2=0.
\end{equation}
This narrows the possibilities to
\begin{equation}
D(0;0,0)
\qquad \text{or} \qquad
D(4;0,0).
\end{equation}

This is the honest large-$R$ plus Casimir conclusion.  It is already nontrivial, because it says that
once the fluctuation is organized globally, the local one-parameter family collapses to two candidate
global modules.  But it still does not choose between them.

\section{Why This Pins Down \texorpdfstring{$\omega$}{omega} for the Primary}

The Casimir by itself labels the whole representation.  It does not tell us the energy of an arbitrary
descendant state inside that representation.  What turns it into a statement about $\omega$ is the
additional fact that we are talking about the \emph{primary}.

Let $H=i\partial_\tau$ be the global AdS Hamiltonian.  A scalar primary state $|\Omega\rangle$ satisfies
\begin{equation}
K_\mu |\Omega\rangle =0,
\qquad
J_{\mu\nu}|\Omega\rangle =0,
\qquad
H|\Omega\rangle = \omega |\Omega\rangle,
\end{equation}
where $K_\mu$ are the conformal-lowering generators and $J_{\mu\nu}$ are the Lorentz generators.
On such a state, the quadratic Casimir reduces to
\begin{equation}
\cC_2 |\Omega\rangle
=
\omega(\omega-4)|\Omega\rangle.
\label{primarycasimir}
\end{equation}
This is just the primary-state version of \eqref{casimir}, with
\begin{equation}
\omega=\D
\end{equation}
for the lowest-weight state.

Combining \eqref{dbranches} and \eqref{primarycasimir}, one gets
\begin{equation}
\omega(\omega-4)=0.
\end{equation}
So the scalar primary can only have
\begin{equation}
\omega=0
\qquad \text{or} \qquad
\omega=4.
\end{equation}

\section{Why the Physical Answer Is \texorpdfstring{$\omega=4$}{omega=4}}

The final step is not algebraic but physical.  The two branches correspond to the two usual AdS scalar
branches:
\begin{equation}
\D=0
\quad\leftrightarrow\quad
\text{source / non-normalizable branch},
\end{equation}
\begin{equation}
\D=4
\quad\leftrightarrow\quad
\text{state / normalizable branch}.
\end{equation}
So the physical axion fluctuation must sit in the normalizable module
\begin{equation}
D(4;0,0),
\end{equation}
and therefore
\begin{equation}
\omega=4.
\end{equation}

\section{What the Derivation Really Uses}

The derivation is conceptually short:
\begin{equation}
\Box_{AdS}\phi=0
\quad + \quad
\text{global } SO(2,4)\text{ covariance}
\quad \Longrightarrow \quad
\cC_2=0.
\end{equation}
Then
\begin{equation}
\cC_2\big(D(\D;0,0)\big)=\D(\D-4)
\quad \Longrightarrow \quad
\D=0 \text{ or } 4.
\end{equation}
Finally,
\begin{equation}
\text{primary} \Rightarrow \omega=\D,
\qquad
\text{normalizable} \Rightarrow \D=4,
\end{equation}
hence
\begin{equation}
\omega=4.
\end{equation}

So the role of the Casimir is not to quantize $\omega$ inside a local patch.  Its role is to identify the
allowed global $SO(2,4)$ module.  Once one knows one is looking at the scalar primary of that module, the
module label immediately becomes the global frequency.

\paragraph{This is exactly where the extra physical input enters.}
If one insists on using only the strict large-$R$ patch analysis plus the representation-theory technology
of \texttt{2504.12290}, then the derivation stops at
\begin{equation}
\cC_2=0
\qquad \Longrightarrow \qquad
D(0;0,0)\ \text{or}\ D(4;0,0).
\end{equation}
To continue from there to
\begin{equation}
\omega=4
\end{equation}
one must still input that the fluctuation is the \emph{normalizable} scalar primary, not the source
branch.  In the compact example of \texttt{2504.12290} the analogue of this extra input is automatic
because the compact-space harmonics are globally regular and the allowed representations are discrete from
the start.  In AdS the group is noncompact, so selecting the physical module requires an additional global
state condition.

\section{Why This Matters for the Current SFT Problem}

This explains exactly what is missing from the large-$R$ patch analysis.  The local radial equation
\eqref{radialode} gives the inner solution, but not the global representation-theoretic choice of module.
The `2504.12290' perspective says that the missing information should come from the globally covariant
fluctuation problem, not from pushing the local series expansion harder.

For the axion this is enough to recover the known answer.  For Konishi one would need the analogue of the
same argument for a coupled higher-spin/stringy fluctuation system rather than a single scalar primary.

\section{Why the Casimir Is Not the Whole Story}

It is important not to overstate the role of the Casimir.  The Casimir does \emph{not} by itself solve the
fluctuation problem.  One still needs genuine dynamical input from the linearized equations of motion.

For the axion, that dynamical input enters in three separate ways:
\begin{itemize}
\item the linearized EOM identify the relevant fluctuation as a massless scalar on AdS$_5$ with trivial
      $S^5$ dependence, so that there are no extra spin, gauge, or mixing terms;
\item the same EOM imply that the geometric operator to match against representation theory is precisely
      $\Box_{AdS}$, with eigenvalue $m_5^2=0$;
\item the local solution of the EOM determines which branch is regular at the center and what the inner
      radial behavior looks like before any global representation-theory condition is imposed.
\end{itemize}

Only after this dynamical reduction does the Casimir become useful.  At that stage it answers a narrower
question:
\begin{equation}
\text{given this fluctuation sector, which global } SO(2,4)\text{ module can it belong to?}
\end{equation}
So the logic is
\begin{equation}
\text{linearized EOM}
\quad \Longrightarrow \quad
\text{effective AdS operator and fluctuation sector},
\end{equation}
\begin{equation}
\text{Casimir matching}
\quad \Longrightarrow \quad
\text{allowed } SO(2,4)\text{ representations},
\end{equation}
\begin{equation}
\text{normalizability / primary condition}
\quad \Longrightarrow \quad
\text{physical } \omega.
\end{equation}

This distinction becomes even more important for Konishi-type fluctuations.  There one should expect:
\begin{itemize}
\item several coupled components,
\item nontrivial tensor or spinor structure,
\item possible constraints or gauge conditions,
\item and mixing between different local field components.
\end{itemize}
In such a case the linearized EOM are needed first to determine the actual reduced system and the operator
whose eigenvalues should be compared to Casimirs.  The Casimir then constrains the global representation
content of that reduced system, but it does not replace the dynamical derivation of the system itself.

Put differently, for the axion the Casimir step is almost trivial \emph{because the dynamics has already
done nearly all of the work}.  Once the linearized EOM reduce the problem to a single massless scalar on
AdS$_5$, the Casimir condition merely says that the fluctuation must belong to one of the two scalar
modules with
\begin{equation}
\cC_2=0 \quad \Longrightarrow \quad \D=0 \ \text{or}\ 4,
\end{equation}
and normalizability then picks the physical branch.  For Konishi, by contrast, one should not expect to
know in advance which combination of the many coupled first-massive components realizes a given
$SO(2,4)\times SO(6)$ module.  There the real problem is first to derive and diagonalize the coupled
linearized system.  Only after that diagonalization can the Casimir be used to identify which irreducible
modules the resulting eigenmodes belong to and how their global energies should be interpreted.

\section{Schematic Konishi Computation}

The Konishi case should therefore be viewed as a problem in degenerate perturbation theory on the full
first-massive string subspace.  Schematically one wants an eigenvalue equation
\begin{equation}
\widehat H_{\psi}\,\Phi=\D\,\Phi,
\end{equation}
for the background-deformed Hamiltonian acting on first-level fluctuations.  Expanding the background and
the effective Hamiltonian in powers of $1/R$,
\begin{equation}
\psi=\frac{1}{R}\psi^{(1)}+\frac{1}{R^2}\psi^{(2)}+\cdots,
\qquad
\widehat H_{\psi}
=
R\,\widehat H^{(0)}+\widehat H^{(1)}+\frac{1}{R}\widehat H^{(2)}+\cdots,
\end{equation}
and expanding the energy as
\begin{equation}
\D=2R+b_0+\frac{b_1}{R}+\cdots,
\end{equation}
one should write the fluctuation as
\begin{equation}
\Phi=\sum_i c_i(R)\,\Phi_i,
\end{equation}
where the $\Phi_i$ form a basis of physical first-massive string fields, ideally organized by
$SO(2,4)\times SO(6)$ quantum numbers.

At zeroth order one has
\begin{equation}
\widehat H^{(0)}\Phi_i=2\,\Phi_i
\end{equation}
on the first-massive subspace, which gives the universal leading term $2R$.

The first nontrivial step is then to restrict $\widehat H^{(1)}$ to this degenerate subspace and
diagonalize the finite matrix
\begin{equation}
M^{(1)}_{ij}=\langle \Phi_i,\widehat H^{(1)}\Phi_j\rangle.
\end{equation}
Its eigenvalues are the candidates for the constant shift $b_0$.  After choosing an eigenvector, one then
computes the next correction from
\begin{equation}
M^{(2)}_{ij}=\langle \Phi_i,\widehat H^{(2)}\Phi_j\rangle
\end{equation}
together with the usual second-order mixing through states outside the degenerate block; this determines
$b_1$.

So the real schematic problem is
\begin{equation}
\det\!\Big[(2R+b_0+b_1/R+\cdots)\mathbf{1}
-
\big(RM^{(0)}+M^{(1)}+R^{-1}M^{(2)}+\cdots\big)\Big]=0,
\end{equation}
with
\begin{equation}
M^{(0)}=2\,\mathbf{1}
\end{equation}
on the first-massive subspace.

This is the point where the contrast with the axion is sharpest.  For the axion, the linearized EOM have
already isolated a single scalar fluctuation, so the Casimir immediately classifies the answer.  For
Konishi, by contrast, the difficult part is first to determine the correct reduced coupled system:
\begin{itemize}
\item include all first-level components that mix under the linearized EOM,
\item quotient pure-gauge and null states,
\item diagonalize block by block,
\item and only then identify the resulting eigenmodes by their global
      $SO(2,4)\times SO(6)$ representations.
\end{itemize}
Only after this dynamical diagonalization should one expect the Casimir to become a clean diagnostic of
which Konishi multiplet state has been found.  For an eigenmode belonging to the Konishi multiplet with
bare label $\D_0$, the expected answer is
\begin{equation}
\D=2R+(\D_0-4)+\frac{2}{R}+\cdots.
\end{equation}

\appendix

\section{Recap of the Logic of \texorpdfstring{\texttt{2504.12290}}{2504.12290}}

This appendix summarizes only the part of \texttt{2504.12290} that is relevant for the axion discussion.
The goal is not to reproduce that paper in detail, but to isolate the logical chain that leads from a
local weak-curvature expansion to a global quantization condition.

\subsection*{A. One local patch is not enough}

The starting point is that an ordinary perturbative SFT expansion around one chosen coordinate origin sees
only a local neighborhood of that point.  This is enough to derive local equations of motion order by
order in $x/R$, but it is not enough to recover global information such as:
\begin{itemize}
\item topology of the target space,
\item global action of the isometry group,
\item or quantization of KK data.
\end{itemize}

\subsection*{B. Introduce a family of normal-coordinate patches}

The key move of \texttt{2504.12290} is to consider not one local coordinate system, but a family of
Riemann normal coordinate systems centered at different target-space points.  The corresponding string
fields are related by diffeomorphism gauge transformations.  Instead of treating these coordinate changes
externally, the paper incorporates them directly into the field-theoretic formalism.

\subsection*{C. Extended BV system}

The ordinary background equation and the transport equations relating nearby coordinate origins are unified
into an extended BV system.  In effect, the extended string field carries dependence on the base point of
the normal-coordinate patch as well as on the worldsheet/string degrees of freedom.

Conceptually, this system does three jobs at once:
\begin{itemize}
\item it imposes the ordinary SFT equations of motion,
\item it implements covariance under infinitesimal diffeomorphisms,
\item and it enforces the integrability conditions required to transport the solution over the target
      manifold.
\end{itemize}

\subsection*{D. Fluctuations as global isometry representations}

Once the background is described globally in this way, fluctuations are no longer just local fields in one
patch.  They can be required to transform covariantly under the global isometry algebra.  This is the step
used in Section 4.4 of \texttt{2504.12290}: the fluctuation equations are supplemented by representation
conditions under the isometry generators.

Applying the representation condition twice produces the quadratic Casimir.  On the geometric side, the
same operator is identified with the target-space Laplacian (plus possible spin terms).  Therefore the
fluctuation spectrum can be read off from representation theory:
\begin{equation}
\text{fluctuation equation}
\quad + \quad
\text{isometry covariance}
\quad \Longrightarrow \quad
\text{Casimir eigenvalue problem}.
\end{equation}

\subsection*{E. Why this matters for the axion}

For the compact example in \texttt{2504.12290}, this logic recovers the expected discrete KK spectrum.  In
the present AdS$_5$ axion problem, the same logic should not be expected to quantize a compact angular
momentum, but rather to select the correct global $SO(2,4)$ module.  The analogue of KK quantization is
therefore the choice of normalizable lowest-weight representation.

That is precisely the missing ingredient in the local projective-patch analysis.  The local radial
equation determines the inner solution, but the global extended/isometry-covariant formulation is what
knows which representation the fluctuation belongs to.  Once that representation is identified as the
normalizable scalar primary module $D(4;0,0)$, the primary frequency is fixed to be
\begin{equation}
\omega=4.
\end{equation}

\end{document}
